\documentclass[aspectratio=169]{beamer}
\input{header}
\coursecode{JKSSB}

\title{OS PYQ Traps \& System Utilities}
\subtitle{Lesson 52 --- Resolving the Classic OS Traps}
\author{JKSSB Computer Awareness}
\date{\today}

\begin{document}
\frame{\titlepage}

\begin{frame}{This Lesson Is a Trap Drill}
  \begin{itemize}
    \item The computer section of JKSSB papers reuses a small set of
      \key{near-neighbour} terms.
    \item A wrong option is usually not random --- it is the \key{closest
      confusion} to the right answer.
    \item In this lesson every question is built around one of those real
      trap patterns.
    \item \key{Read each option as a claim to be tested}, not merely as a
      phrase to recognise.
    \item \muted{Keep one fact per line and preserve the exact English term.}
  \end{itemize}
  \begin{alertblock}{Trap drill}
    The wrong option is usually the ``sounds right'' neighbour --- check the
    exact wording before you choose.
  \end{alertblock}
\end{frame}

\begin{frame}{The Operating System: What It Is}
  \begin{itemize}
    \item The \key{operating system (OS)} is \key{system software}.
    \item It provides a \key{user-friendly interface} between the user and
      the hardware (PYQ-21).
    \item It manages hardware resources: processor, memory, storage, and
      input/output devices.
    \item It provides services to application programs.
    \item \muted{Examples: Windows, Linux, UNIX, MS-DOS, Android, macOS.}
  \end{itemize}
  \begin{block}{Definition}
    The OS is the system software that controls the hardware and offers a
    user interface and services to programs.
  \end{block}
\end{frame}

\begin{frame}{Where the OS Lives While Running}
  \begin{itemize}
    \item An OS is stored permanently on \key{secondary storage} (the disk).
    \item While the computer is \key{running}, the OS is loaded into
      \key{primary memory (RAM)} (PYQ-22).
    \item The CPU can execute OS code only from main memory.
    \item \muted{Power off clears RAM; the OS is reloaded from disk at the
      next boot.}
  \end{itemize}
  \begin{alertblock}{Trap alert}
    The OS is not executed from secondary storage. At run time it resides in
    primary memory.
  \end{alertblock}
\end{frame}

\begin{frame}{Three Software Categories}
  \begin{center}
    \begin{tabular}{p{0.24\textwidth}p{0.34\textwidth}p{0.32\textwidth}}
      \toprule
      \textbf{Category} & \textbf{Role} & \textbf{Examples} \\
      \midrule
      System software & Runs and manages the machine & OS, device drivers, utility programs, translators \\
      Application software & Does a user's task & MS Word, Excel, a browser \\
      Utility software & Maintains and optimises the system & Disk Cleanup, Disk Defragmenter, antivirus \\
      \bottomrule
    \end{tabular}
  \end{center}
  \vspace{0.25cm}
  \muted{Utility software is a maintenance category closely tied to system
  software.}
\end{frame}

\begin{frame}{Device Drivers Are Software}
  \begin{itemize}
    \item A \key{device driver} is a small \key{software program} that lets
      the OS talk to a specific hardware device.
    \item It translates OS commands into device-specific instructions.
    \item Every device --- printer, keyboard, graphics card --- needs a
      matching driver.
    \item \muted{The driver is software; the device it controls is
      hardware.}
  \end{itemize}
  \begin{alertblock}{Trap alert}
    A device driver is \key{software}, not hardware (TRAP-10). The OS itself
    is \key{system software}.
  \end{alertblock}
\end{frame}

\begin{frame}{DMA: Transferring Data Without the CPU}
  \begin{itemize}
    \item \key{DMA} stands for \key{Direct Memory Access}.
    \item A DMA controller transfers data directly between an \key{I/O
      device} and \key{main memory}.
    \item The CPU does \key{not} handle each individual byte of the transfer
      --- it is bypassed.
    \item The CPU starts the transfer and is notified (by interrupt) when it
      finishes.
    \item This frees the CPU for other work and speeds up bulk transfers.
  \end{itemize}
  \begin{block}{Definition}
    DMA is a method in which a controller moves data between a device and
    main memory without the CPU's continuous involvement.
  \end{block}
\end{frame}

\begin{frame}{DMA Trap Alert}
  \begin{itemize}
    \item DMA moves data between \key{an I/O device and main memory}, not
      between the CPU and the device.
    \item The CPU is \key{bypassed} during the transfer, not made faster.
    \item DMA does \key{not} increase RAM size and does not encrypt data.
    \item \muted{The CPU still initiates and completes the transfer.}
  \end{itemize}
  \begin{alertblock}{Trap alert}
    ``DMA = the CPU does the copying faster'' is wrong. DMA removes the CPU
    from the per-byte copying.
  \end{alertblock}
\end{frame}

\begin{frame}{RTOS: Real-Time Operating System}
  \begin{itemize}
    \item \key{RTOS} stands for \key{Real-Time Operating System}.
    \item It is designed so that tasks finish within \key{fixed time
      deadlines}.
    \item It is used in embedded and control systems: machinery, medical
      devices, vehicle electronics.
    \item \muted{Chosen for timing guarantees, not for word processing or
      graphics design.}
  \end{itemize}
  \begin{block}{Definition}
    An RTOS is an operating system that guarantees a response within a
    bounded time.
  \end{block}
\end{frame}

\begin{frame}{DOS, UNIX, and Linux Distributions}
  \begin{itemize}
    \item \key{DOS} = \key{Disk Operating System}; a command-line OS
      (MS-DOS).
    \item \key{UNIX} is a multi-user, multitasking operating system.
    \item \key{Linux} is an open-source, UNIX-like OS; it comes in
      \key{distributions}.
    \item Common Linux distributions: \key{Ubuntu}, \key{Fedora},
      \key{Debian}, \key{Linux Mint}.
    \item \muted{Windows and macOS are operating systems too, but they are
      not Linux distributions.}
  \end{itemize}
\end{frame}

\begin{frame}{Exact OS Recognition}
  \begin{itemize}
    \item Ask first: does the option actually name an \key{operating
      system}?
    \item Operating systems: \key{Windows}, \key{Linux}, \key{UNIX},
      \key{MS-DOS}, \key{Android}, \key{macOS}.
    \item \key{Not} operating systems: MS Word, MS Excel, MS Access,
      Chrome, Firefox --- these are applications.
    \item \muted{MS Access is a database application, not an operating
      system.}
  \end{itemize}
  \begin{alertblock}{Trap alert}
    A familiar brand name alone does not make a system an OS --- check
    whether it is system software.
  \end{alertblock}
\end{frame}

\begin{frame}{Internal vs External Hardware}
  \begin{center}
    \begin{tabular}{p{0.22\textwidth}p{0.30\textwidth}p{0.38\textwidth}}
      \toprule
      \textbf{Location} & \textbf{Meaning} & \textbf{Examples} \\
      \midrule
      Internal & Inside the system unit / cabinet & motherboard, CPU, RAM, hard disk \\
      External (peripheral) & Outside the system unit, joined by ports & monitor, keyboard, mouse, printer, scanner \\
      \bottomrule
    \end{tabular}
  \end{center}
  \vspace{0.25cm}
  \muted{The classification depends on position relative to the system unit.}
\end{frame}

\begin{frame}{The Windows Run Box and \texttt{winword}}
  \begin{itemize}
    \item The Windows \key{Run} box (Windows + R) starts a program when you
      type its command.
    \item Typing \key{\texttt{winword}} opens \key{MS Word}.
    \item \key{\texttt{notepad}} opens Notepad; \key{\texttt{wordpad}} opens
      WordPad.
    \item \key{\texttt{cmd}} opens the Command Prompt.
    \item \muted{Match the exact command to the exact program.}
  \end{itemize}
  \begin{alertblock}{Trap alert}
    \texttt{winword} = MS Word. Do not confuse it with \texttt{wordpad} (a
    different editor) or \texttt{notepad}.
  \end{alertblock}
\end{frame}

\begin{frame}{Trap Alert: Terms to Keep Straight}
  \begin{itemize}
    \item Device driver = \key{software}; OS = \key{system software}
      (TRAP-10).
    \item DMA moves data between a device and main memory, \key{bypassing
      the CPU}.
    \item RTOS = \key{Real-Time Operating System} (deadline-driven).
    \item Utility = maintenance software; application = the user's task.
    \item Linux distributions: Ubuntu, Fedora, Debian; \key{MS Access is
      not an OS}.
    \item \texttt{winword} opens MS Word.
  \end{itemize}
\end{frame}

\begin{frame}{Lesson 52: Recall}
  \begin{itemize}
    \item The OS is system software; it provides a user interface, manages
      resources, and runs from primary memory.
    \item Software categories: system, application, and utility.
    \item A device driver is \key{software}; the OS is \key{system
      software}.
    \item DMA = Direct Memory Access; data moves between device and main
      memory with the CPU bypassed.
    \item RTOS = Real-Time Operating System.
    \item DOS = Disk Operating System; Linux distributions include Ubuntu,
      Fedora, and Debian.
    \item Internal vs external hardware; \texttt{winword} opens MS Word.
  \end{itemize}
\end{frame}
\end{document}
